% Figure: steady-state amplitude versus drive frequency for a hardening
% Duffing oscillator. The response curve bends to the right and folds over,
% so a swept drive jumps down on the way up and jumps up on the way down,
% tracing a hysteresis loop between two turning points.
\begin{figure}[htbp]
\centering
\begin{tikzpicture}
\begin{axis}[
  width=12cm, height=8cm,
  xlabel={drive frequency $\Omega$}, ylabel={steady-state amplitude $A$},
  xmin=0.6, xmax=1.9, ymin=0, ymax=3.4,
  axis lines=left,
  grid=major, grid style={enggray!20},
  tick label style={font=\scriptsize},
  label style={font=\small},
  legend pos=north west,
  legend cell align=left,
  legend style={font=\scriptsize},
]
% Upper (large-amplitude) stable branch: bends right, ends at upper turning point
\addplot[engblue, very thick] coordinates {
  (0.62,0.35) (0.75,0.5) (0.9,0.75) (1.05,1.1) (1.2,1.55)
  (1.35,2.1) (1.5,2.7) (1.62,3.15) (1.66,3.3)
};
\addlegendentry{stable (upper branch)}
% Lower stable branch
\addplot[engorange, very thick] coordinates {
  (0.62,0.32) (0.8,0.42) (1.0,0.55) (1.15,0.66) (1.28,0.78)
  (1.4,0.9) (1.5,1.02) (1.6,1.16) (1.7,1.32) (1.72,1.38)
};
\addlegendentry{stable (lower branch)}
% Unstable middle branch (dashed) connecting the two turning points
\addplot[engred, thick, dashed] coordinates {
  (1.66,3.3) (1.6,2.55) (1.5,1.95) (1.4,1.55) (1.32,1.35)
  (1.25,1.2) (1.72,1.38)
};
\addlegendentry{unstable branch}
% jump-down arrow at the upper turning point (sweep up)
\draw[engblue, -{Latex[length=2mm]}, thick] (axis cs:1.66,3.15) -- (axis cs:1.68,1.42);
\node[engblue, font=\scriptsize, anchor=west] at (axis cs:1.68,2.3) {jump down (sweep up)};
% jump-up arrow at the lower turning point (sweep down)
\draw[engorange, -{Latex[length=2mm]}, thick] (axis cs:1.25,1.2) -- (axis cs:1.23,1.9);
\node[engorange, font=\scriptsize, anchor=east] at (axis cs:1.22,0.9) {jump up (sweep down)};
\end{axis}
\end{tikzpicture}
\caption[Frequency response of a hardening Duffing oscillator]{Steady-state
amplitude against drive frequency for a hardening Duffing oscillator,
$\ddot{u}+2\zeta\dot{u}+u+\varepsilon u^{3}=F\cos\Omega t$ with
$\varepsilon>0$. The stiffening cubic term pulls the resonance peak to
the right and folds the response curve over, so that across a band of
frequencies three steady states coexist: two stable (blue and orange
solid) and one unstable (red dashed) that no physical run ever settles
on. A drive swept slowly upward in frequency rides the upper branch to
the right-hand turning point and then jumps down to the lower branch;
a drive swept downward rides the lower branch to the left-hand turning
point and jumps up. The two sweeps trace different paths, the
hysteresis loop that a linear resonance never shows. The single
crossing of the linear case has become a fold, and the amplitude the
system displays now depends on the history of the sweep.}
\label{fig:vol02:ch07:duffing-response}
\end{figure}
